\documentclass[12pt]{article}
\usepackage{geometry}
\usepackage{amsfonts}
\usepackage{float}
\usepackage{amsmath}

\geometry{
legalpaper, total={177.8mm, 290mm},left=5mm, right=15mm,
top=7mm, bottom=12mm,
}

\usepackage{letltxmacro}

\LetLtxMacro\Oldfootnote\footnote

\newcommand{\EnableFootNotes}{%
  \LetLtxMacro\footnote\Oldfootnote%
}

\newcommand{\DisableFootNotes}{%
  \renewcommand{\footnote}[2][]{\relax}
}

\DisableFootNotes

\begin{document}

\iffalse
\begin{table}[]
\begin{tabular}{llcllcll}
 &  & \textbf{MYMENSINGH GIRLS’  CADET COLLEGE} &  &                                            & \multicolumn{1}{l}{}            &                        &                        \\
 &  & WHAT EXAMINATION - 2025                   &  &                                            &                                 &                        &                        \\
 &  & CLASS: WHAT                               &  &                                            &                                 &                        &                        \\
 &  & STATISTICS (CREATIVE)                     &  &                                            &                                 &                        &                        \\
 &  & WHAT PAPER                                &  &                                            & \multicolumn{1}{r}{}            &                        &                        \\ \cline{7-8} 
 &  & [According to the Syllabus of 2025]       &  & \multicolumn{1}{r}{\textbf{Subject Code:}} & \multicolumn{1}{l|}{\textbf{1}} & \multicolumn{1}{l|}{3} & \multicolumn{1}{l|}{0} \\ \cline{7-8} 
 &  & TIME – 2 hours \& 35 minutes              &  &                                            & \multicolumn{1}{r}{}            &                        &                        \\
 &  & FULL MARKS – 50                           &  &                                            & \multicolumn{1}{r}{\textbf{}}   &                        &                       
\end{tabular}
\end{table}
\fi

\begin{table}[]
\begin{tabular}{llllllllcllcl}
\textit{Elnath} &  &  &  &  &  &  &  & \textbf{MYMENSINGH GIRLS’ CADET COLLEGE} &                                             &                                 & \multicolumn{1}{l}{}            &                        \\
                 &  &  &  &  &  &  &  & PROGRESS TEST EXAMINATION - 2026          &                                             &                                 &                                 &                        \\
                 &  &  &  &  &  &  &  & CLASS: HSC                               &                                             &                                 &                                 &                        \\
                 &  &  &  &  &  &  &  & STATISTICS (CREATIVE)                    &                                             &                                 &                                 &                        \\
                 &  &  &  &  &  &  &  & SECOND PAPER                             &                                             &                                 & \multicolumn{1}{r}{}            &                        \\ \cline{11-13} 
                 &  &  &  &  &  &  &  & [According to the Syllabus of 2026]      & \multicolumn{1}{r|}{\textbf{Subject Code:}} & \multicolumn{1}{l|}{\textbf{1}} & \multicolumn{1}{l|}{\textbf{3}} & \multicolumn{1}{l|}{0} \\ \cline{11-13} 
                 &  &  &  &  &  &  &  & TIME – 1 hours \& 25 minutes             &                                             &                                 & \multicolumn{1}{r}{}            &                        \\
                 &  &  &  &  &  &  &  & FULL MARKS – 30                          &                                             &                                 & \multicolumn{1}{r}{\textbf{}}   &                       
\end{tabular}
\end{table}


\hrule

\begin{center}
[\textbf{N.B.} – The figures of the right margin indicate full marks. Read the stems carefully and answer the associated questions. Answer any \textbf{THREE} questions taking at least ONE question from each group]\\
\end{center}

\begin{center}
\textbf{Group  - A}
\end{center}
 \begin{enumerate}
 
 \item
  \textbf{In a class of 40 students, 23 passed in Mathematics, 25 passed in English, and 10 failed in both the subjects. A student is randomly selected.}

  \begin{enumerate}
    \item
	What does $P(A) = 0$ mean? \hfill 1
    \item
	Find the probability that the student passed in at least one subject. \hfill 2
    \item  
	What is the probability that the student passed in both the subjects? \hfill 3
    \item
	Determine the probability that the student passed in only one subject. \hfill 4
  \end{enumerate} 

\item
  \textbf{The monthly electricity bills (in dollars) of 5 households are 45, 60, 55, 70, and 50. A researcher claims that if random samples of size 3 are drawn from these bills without replacement, the sample mean will provide an unbiased estimate of the population mean.}

  
  \begin{enumerate}
    \item
    Define a sample in this context. \hfill 1
    \item
    State two drawbacks of conducting a census. \hfill 2
    \item  
    Compute the population variance from the given data. \hfill 3
    \item
    Verify the researcher's claim regarding unbiasedness of the sample mean. \hfill 4
  \end{enumerate}




  \begin{center}
\textbf{Group  - B}
\end{center}

  \item
  \textbf{The joint probability function of two random variables \( X \) and \( Y \) is given by:}
  
  \begin{center}
  \( \displaystyle P(X,Y) = \frac{x + y + 1}{42}; \quad x = 0, 1, 2; \quad y = 0, 1, 2, 3 \)
  \end{center}
 
  \begin{enumerate}
    \item What is a discrete random variable? \hfill 1
    \item What are the required properties of a probability distribution? \hfill 2
  
    \item
    	Calculate the marginal probability \( P(Y) \). \hfill 3
    \item
     	Determine \( P(Y \vert X = 1) \) and \( P(Y \vert X = 0) \). \hfill 4
  \end{enumerate}
  
  
 \item
  \textbf{The population of a country is growing exponentially at an annual rate of 2.5\%.}

  
  \begin{enumerate}
    \item
	How does the arithmetic growth increase? \hfill 1
    \item
    What is the key difference between exponential growth and geometric growth? \hfill 2
  \item
How many years will it take for the population to double under geometric growth and \\ exponential growth? \hfill 3

\item
If the current population is 8 million, after 56 years, how many **additional** people will \\ there be if the population grows exponentially instead of arithmetically? \hfill 4
  \end{enumerate}

\end{enumerate}

 \vspace{2.5cm}

\begin{center}

  \vfill
  --AAM--
\end{center}



\end{document}